\chapter{Arithmetic with Memory: the Semiring}
\label{ch:arithmetik}

% Register: D6, S5, S6, S10-S16, V9. Sources: Spec. 8, Appendices A-C.

Up to this point we have only looked at Hori numbers and shifted
them about. Now we want to compute with them. That sounds harmless,
but it is the point at which Horimetrik demands its first genuine
decision: when two Hori numbers are added, which horizon does the
result receive?

In the ordinary number system this question never arises -- there a
result has no horizon, only a value. Here, by contrast, the horizon
is part of the number. A rule of computation must therefore consist
of two parts: a rule for the value and a rule for the horizon. The
horizon rule we call a \emph{selector}.

\section{The first failure and what it teaches}

The most obvious selector is the conservative one: the result
inherits the largest horizon involved, enlarged if necessary to
what the result value minimally requires. For addition this works
flawlessly. For multiplication it fails -- and it fails at the most
innocent of all numbers.

\begin{sackgasse}
If one multiplies conservatively, then in general
$(a\htimes b)\htimes c\ne a\htimes(b\htimes c)$ as soon as a zero
is involved. The zero erases the \emph{value} of every product, but
the conservative selector keeps dragging along the \emph{horizons}
of the intermediate steps. Depending on when the zero strikes, the
result remembers different computational routes -- same value,
different horizons, no freedom of bracketing. The lesson: a horizon
carries computational history, and computational history does not
automatically get along with the laws of arithmetic.
\end{sackgasse}

\begin{beispiel}
Concretely, with the conservative multiplication (horizon = maximum
of those involved, enlarged if necessary): let $a=b=(2,5)$ -- the
five with horizon 2 -- and $c=(1,0)$ the zero with horizon 1.
Bracketed on the left: $a\htimes b=(4,25)$, since $25$ needs
horizon 4; then $(4,25)\htimes c=(4,0)$. Bracketed on the right:
$b\htimes c=(2,0)$, then $a\htimes(2,0)=(2,0)$. Both routes end at
the value zero -- but once with horizon 4, once with horizon 2.
$(4,0)\ne(2,0)$: no associativity.
\end{beispiel}

\section{The value-side semiring}

The way out is a division of labour: addition may remember,
multiplication must forget. (One convention in advance, valid for
the whole book: we call an arithmetic \emph{law-abiding} if both
operations are associative and commutative, multiplication
distributes over addition, and an additively neutral element
exists. A multiplicative identity and absorption by the additive
zero we do \emph{not} require in general, but only where they are
explicitly named -- both are about to meet us as genuine
distinguishing features.)

\begin{satz}[value-side semiring]\label{satz:a3}
On the Hori space $\Hori=\{(n,x)\mid n\ge\muh(x)\}$ let, in value
coordinates $(n,x)$,
\[
(n,x)\hplus(m,y)=\bigl(\max(n,m,\muh(x+y)),\,x+y\bigr),
\qquad
(n,x)\htimes(m,y)=\bigl(\muh(xy),\,xy\bigr).
\]
Then $\hplus$ and $\htimes$ are commutative and associative,
$\htimes$ distributes over $\hplus$, and $(0,0)$ is neutral for
$\hplus$ and absorbing for $\htimes$.
\end{satz}

The proof is astonishingly short; its single core is the
monotonicity of the minimal horizon: more value never needs less
room (in full in Appendix~\ref{ch:anhang-beweise}).
What is remarkable is what this semiring \emph{lacks}: an identity.
There is no element that leaves everything unchanged under
multiplication. Instead, for the Hori number one -- that is,
$(1,1)$, the one in its smallest home -- we have the equation
\[
h\htimes(1,1)=\KV(h)
\qquad\text{with}\quad
\KV(n,x)=(\muh(x),\,x).
\]
Multiplication by one is the value
compactification.\footnote{Terminological warning:
``compactification'' is not meant topologically here. Formally
these are idempotent retractions onto the respective minimal
horizon -- we keep the intuitive word and ask topologically trained
readers for their indulgence.} The one acts
not as a neutral element but as a \emph{projector}: it presses
every Hori number down to its compact form. The compact elements
form a small semiring of their own underneath, and it is exactly
the ordinary $\Nn$. The familiar natural numbers thus sit as a
``corner'' in the Hori space, and the way there is not an extra
operation but multiplication by one.

\begin{beispiel}
Let us work through it once. $(3,7)\hplus(2,5)$: value $12$,
horizon $\max(3,2,\muh(12))=\max(3,2,3)=3$, hence $(3,12)$ -- the
largest horizon involved survives. $(3,7)\htimes(2,5)$: value $35$,
horizon $\muh(35)=4$, hence $(4,35)$ -- the past is forgotten. And
the projector: the seven with an unnecessarily large horizon,
$(5,7)$, times the one $(1,1)$ gives
$(\muh(7),7)=(3,7)$ -- multiplication by one has tidied up.
\end{beispiel}

\begin{oberflaeche}
What has happened here? We have found an arithmetic in which
numbers remember their past when added -- the largest horizon
involved survives -- and are born anew when multiplied. And the
ordinary arithmetic we have known since primary school is contained
within it: as the part of the Hori space that carries no past.
\end{oberflaeche}

\section{Excess: memory acquires a coordinate}

For a long time it looked as though this forgetting multiplication
had no alternative. The natural conjecture was that the laws of
arithmetic destroy every other horizon rule. It is false, and the
refutation is more instructive than a proof would have been.

The horizon of every Hori number splits canonically into two parts:
\[
n=\underbrace{\muh(x)}_{\text{obligation}}
 +\underbrace{\varepsilon}_{\text{excess}} .
\]
The obligatory part is the minimum the value needs; the excess is
memory carried along voluntarily. This decomposition is not a
convention but is uniquely determined by the factorial shells --
and it decouples the question of the selector completely: one may
compute on the excess \emph{independently}.

\begin{satz}[product principle]\label{satz:produktprinzip}
Every commutative semiring structure on the excesses yields a
law-abiding Hori arithmetic: values compute in the ordinary way,
excesses compute in the chosen structure.
\end{satz}

The product principle thus supplies a whole family of new
arithmetics. Together with the forgetting semiring already
constructed in Theorem~\ref{satz:a3} -- which, mind you, is
\emph{not} a product construction: its addition selector compares
the horizons themselves, not merely the excesses -- three
particularly natural arithmetics emerge:

\begin{center}
\small
\begin{tabular}{lccc}
 & memory under $+$ & memory under $\times$ & identity \\
\hline
forgetting (Thm.~\ref{satz:a3}; not a product) & maximum (of horizons) & dies & none \\
preserving (tropical) & maximum & adds up & $(1,1)$ \\
amplifying (arithmetic) & adds up & multiplies & $(2,1)$ \\
\end{tabular}
\end{center}

\begin{beispiel}
Let $a=(5,7)$ -- obligation $\muh(7)=3$, excess $2$ -- and
$b=(3,5)$ -- obligation $2$, excess $1$. The product has value $35$
with obligation $\muh(35)=4$. The three arithmetics differ only in
the excess of the result: forgetting $(4,35)$ (excess $0$),
preserving $(4{+}2{+}1,35)=(7,35)$ (excesses add up),
amplifying $(4{+}2,35)=(6,35)$ (excesses multiply:
$2\cdot1=2$).
\end{beispiel}

The preserving arithmetic has a curiosity: its zero does not
absorb. If one multiplies a Hori number by zero, the value is wiped
out, but the memory survives. The amplifying arithmetic, by
contrast, is a fully fledged commutative semiring with identity and
absorbing zero -- and its multiplication does not erase memory
wholesale but reckons with it arithmetically. Admittedly it too is
not entirely faithful to memory: an operand with excess zero
absorbs the other's memory, since zero times anything remains
zero -- and exactly this is where its absorbing zero comes from.

\begin{oberflaeche}
The question ``which horizon does the result receive?'' has turned
into something far more beautiful: ``what is to happen to the past
when computation takes place?'' And Horimetrik's answer is: that is
yours to choose -- forget, preserve, or amplify -- and each choice
yields a world that is perfectly free of contradiction within
itself.
\end{oberflaeche}

\section{The duality: the same one, two projectors}

Up to this point everything computed on the value side. But Hori
numbers carry, besides their value, their shape as well (and, since
Chapter~\ref{ch:erweiterungen}, the fraction -- of which more in
Section~\ref{sec:seitentreue}) -- and the shape side has
an arithmetic of its own. Instead of the
values one can add and multiply the \emph{structure polynomials};
that is digit-wise computation without carries, in which, for
instance, the structures of $4$ and $5$ directly give rise to the
structure of $20$.

\begin{satz}[duality theorem]\label{satz:dualitaet}
The Hori space carries a second, structure-side semiring: addition
conservative, multiplication compact -- but measured against the
structure horizon $\sigmah$ instead of the value horizon $\muh$.
Computation takes place in structure coordinates $(n,P)$; via the
digit bijection from Chapter~\ref{ch:horizonte}, $(n,x)$ and the
pair consisting of $n$ and the shape of $x$ denote the same
element. In this semiring
\[
h\boxtimes(1,1)=\KS(h)
\qquad\text{with}\quad
\KS(n,P)=(\sigmah(P),\,P),
\]
and $\KS$ is compatible with both operations. The corner of this
semiring is the space of structure polynomials.
\end{satz}

The punchline deserves a paragraph of its own. It is the
\emph{same} Hori number $(1,1)$ -- the one -- that carries the
projector $\KV$ on the value side and the projector $\KS$ on the
structure side. Two multiplications, one element. And the oldest
result of this project -- that the two tidying operations do not
commute; the next section makes this formal -- turns out to be the
non-commuting of these two multiplications with one and the same
one.

\section{Tidying up: the commutator and the depth law}
\label{sec:kommutator}

Because a whole thematic strand of this book rests on them, we now
place the two projectors side by side (Register D6):
\[
\KV(n,x)=(\muh(x),\,x),
\qquad
\KS(n,P)=(\sigmah(P),\,P).
\]
$\KV$ tidies up \emph{value-compactly}: same value, smallest
horizon -- the shape is recoded in the process. $\KS$ tidies up
\emph{structure-compactly}: same shape, smallest horizon -- the
value migrates downwards with the evaluation point. Both are
idempotent, but they do not commute, and the deviation is
measurable. We call
\[
\delta(h)=\hor\bigl(\KS\KV(h)\bigr)-\hor\bigl(\KV\KS(h)\bigr)
\]
the \emph{compactification commutator}; its negative values are
called \emph{depth}.

\begin{beispiel}
The house spirit of this book: $h=(6,6)$, the six in terminal
representation. $\KV$ first: $(6,6)\mapsto(3,6)$; there the six has
the shape $X{+}2$ with $\sigmah=2$, hence
$\KS\KV(h)=(2,5)$ -- horizon~$2$. $\KS$ first: the constant shape
$6$ is already structure-compact, $\KS(h)=h$, then
$\KV(h)=(3,6)$ -- horizon~$3$. Hence $\delta(h)=2-3=-1$: the order
of tidying up decides where you end up -- and as what.
\end{beispiel}

Upwards, $\delta$ grows comparatively fast; the smallest witness
for $+2$ is the $720$ as shape $\ff X3$ in horizon $9$
(the ``single one in the fourth place'' from the last
interlude). Downwards, by contrast, $\delta$ is factorially
sluggish, and this sluggishness obeys an exact law. Call
\[
g(m)=m-\min\{s\mid M_s(m{+}1)\ge m!\}
\]
the \emph{shell gap} of the shell $S_m$: how far below $m$ the
capacity of a shape may lie whose value nevertheless still reaches
$S_m$ ($M_s$ is the maximally filled shape from
Chapter~\ref{ch:strukturpolynome}).

\begin{satz}[depth law, lower side; Register S33]
\label{satz:tiefengesetz}
For every shell $m\ge3$ we have
\[
\min\{\delta(h)\mid \muh(\val(h))=m\}=-g(m),
\]
and an always serviceable witness -- in general not the only
one -- is the terminal representation of the maximal-polynomial
value $M_{m-g(m)}(m{+}1)$. The jump positions
$m_g=\min\{m\mid g(m)\ge g\}$
form a counting sequence with proven asymptotics
$m_g/(g{+}2)!\to1$ (proofs in Appendix~\ref{ch:anhang-beweise}).
\end{satz}

For shell $3$ the witness is exactly the house spirit from above,
and $m_1=3$: from the third shell onwards depth $-1$ is possible,
from the fifteenth $-2$, from the eighty-fourth $-3$. The
\emph{two-sided fixed points}, finally, are the elements that
tidying up can no longer touch -- fixed by $\KV$
\emph{and} $\KS$ at once, that is, $n=\muh(x)=\sigmah(P)$. Both
notions, jump positions and fixed points alike, are counted in
Chapter~\ref{ch:zaehlungen}.

One fine-grained detail of this non-commutativity deserves mention:
the transformation monoid generated by $\KV$ and $\KS$ -- the
\emph{compactification monoid} -- is finite: exactly six elements,
all of which are idempotent except $\KV\KS$
(Register S21; proof in Appendix~\ref{ch:anhang-beweise}). The
witness with the smallest \emph{value} is again the house spirit
$(6,6)$; in the smallest \emph{world} the idempotence of
$\KV\KS$ already fails at $(5,25)$.

\section{Side fidelity}
\label{sec:seitentreue}

Can the two sides be interwoven -- say, adding on the value side
and multiplying on the structure side? Law-abiding here means, as
agreed at the outset: both operations commutative and associative,
multiplication distributes, an additive zero exists. The systematic
test of all sixteen combinations -- value or shape as payload,
conservative or compact as selector, for each of the two
operations -- then gives a clear answer: exactly two \emph{of these
sixteen} are law-abiding, the value-side semiring and its
structure-side mirror. (The product family from
Theorem~\ref{satz:produktprinzip} stands outside this grid -- its
selectors compute on excesses -- and is law-abiding throughout, but
stays entirely on the value side.) All eight mixed candidates die
of distributivity; the remaining six side-pure ones fail at
neutrality (compact addition forgets the operand's horizon) or at
associativity (conservative multiplication, the zero problem from
the beginning of the chapter).

\begin{beispiel}
The paradigmatic witness, digit by digit. (In this example $\oplus$
and $\otimes$ denote the mixed candidate arithmetic, not the
operations from Theorem~\ref{satz:a3}; \emph{value-carried} means:
the operation reckons with the values, \emph{structure-carried}:
with the shapes.) Let $a=b=(1,1)$ and
$c=(2,5)$, whose shape is $X+2$ (digits $12$). Adding
value-carried, $b\oplus c=(3,6)$; in $H_3$ the six happens to have
the shape $X+2$ as well (digits $012$). Multiplying now
structure-carried by $a$ (shape $1$), the shape $X+2$ persists and
is evaluated compactly at horizon 2: result
$(2,5)$ -- \emph{the value has dropped from 6 to 5}, because the
structure multiplication knows only the shape. On the right, by
contrast: $a\otimes b=(1,1)$, $a\otimes c=(2,5)$, added
value-carried $(3,6)$. Left $(2,5)$, right $(3,6)$: distributivity
is dead, and one sees exactly why -- one side honestly keeps
computing with the value, the other lost it in the change of
shape.
\end{beispiel}

\begin{satz}[side fidelity theorem]\label{satz:seitentreue}
Every law-abiding arithmetic \emph{of the payload-carried form
considered here} -- the operations carry value, shape,
or fraction, the horizon rules are arbitrary -- is side-pure:
addition and multiplication both operate on the value side or
both on the structure side.
\end{satz}

The restriction to payload-carried operations is essential: on a
countable set any semiring structure whatsoever can be defined by
transport -- the theorem speaks about the arithmetics that
\emph{respect} value, shape, or fraction, and about these it
speaks completely.

On its exact reach: depending on the case, the proof needs even
less than full law-abidingness -- one mixed case dies of left
distributivity alone, the other of two-sided distributivity (which
belongs to every semiring, including non-commutative ones): right
distributivity mirrors the chain construction and replaces
commutativity. Non-commutative mixed forms are thereby excluded as
well; all that remains conceivable is a mixed system that is
distributive on \emph{one side} only -- a curiosity outside any
semiring axiomatics.

The proof deserves to have its core idea told, for it has a
punchline (in full in Appendix~\ref{ch:anhang-beweise}). At first things looked bad: individual
computations that a mixed arithmetic would have to satisfy are in
fact solvable -- in part via Pell equations, that is, exactly the
kind of number theory that likes to produce exceptions. A direct
contradiction seemed out of reach. The decisive weapon is
astonishingly plain: the \emph{constants}. If one multiplies a
chain of numbers built up by addition with ever larger constant
structures, the digit bound drives the evaluation horizons to
infinity -- and a non-constant polynomial cannot take the same
value at arbitrarily many points. That forces the chain links to
be constant. Constants, however, die at the next multiplier,
because there too the digit bound makes the horizon grow linearly,
while distributivity demands linear growth of the value. Both mixed
forms break at the very same spot.

By now the theorem even has a third dimension: the fraction
identity from Chapter~\ref{ch:erweiterungen} can be interwoven
neither with the value side nor with the shape side, and it carries
no arithmetic of its own (its addition can exceed
one). Side fidelity is a trilogy: three identities,
two computational worlds, no mixture -- proofs in
Appendix~\ref{ch:anhang-beweise}.

The punchline: the source of the rigidity is the digit condition
$0\le a_i\le i$ -- the very first, most innocent rule of this book.
The rule with which everything began is exactly the rule that
prevents value and shape from ever merging into one arithmetic. It
is not the horizon rules that are rigid; rigid are the identities.

\begin{oberflaeche}
A Hori number has three identities: what it is worth, what it
looks like, and which fraction it is. This chapter has shown that
the first two carry fully fledged, mirror-image computational
worlds, that one and the same one works as doorkeeper in both
worlds -- and, as a proven theorem: that nobody can compute in two
worlds at once, while the third carries none of its own. One must
choose. That is the deepest theorem Horimetrik has to offer:
consistency has a price, and the price is an identity.
\end{oberflaeche}
